% =====================================================================
%  مقدمة الناظم (٤): الأبيات ٦٧–٩٤ — وصف القصيدة وتسميتها، ثم الدعاء والوصايا
% =====================================================================
\chapter{مقدمة الناظم (٤): وصف القصيدة والوصايا}

\section{وصف القصيدة وصلتها بكتاب التيسير}

\bayt{67}{أَهَلَّتْ فَلَبَّتْهَا الْمَعَانِي لُبَابُهَا}{وَصُغْتُ بِهَا مَا سَاغَ عَذْبًا مُسَلْسَلَا}
\begin{mufradat}
\kalima{أهلّت}{رفعت صوتها، أي نادت.}
\kalima{لبّتها}{أجابتها.}
\kalima{لُبابها}{خالصها وأحسنها.}
\kalima{صُغت}{أحكمت النظم (من الصياغة).}
\kalima{ساغ}{سهُل وطاب.}
\kalima{مسلسلًا}{سلِسًا صافيًا.}
\end{mufradat}
\maana نادت القصيدةُ المعانيَ فأجابها خيارها، ونظمتُ فيها لفظًا سهلًا عذبًا يجري على اللسان.

\bayt{68}{وَفِي يُسْرِهَا التَّيْسِيرَ رُمْتُ اخْتِصَارَهُ}{فَأَجْنَتْ بِعَوْنِ اللهِ مِنْهُ مُؤَمَّلَا}
\begin{mufradat}
\kalima{التيسير}{كتاب «التيسير في القراءات السبع» للإمام أبي عمرو الداني (ت ٤٤٤ هـ).}
\kalima{رُمتُ}{قصدتُ.}
\kalima{أجنت}{أثمرت وكثر جناها.}
\kalima{مؤمّلًا}{مرجوًّا.}
\end{mufradat}
\maana قصدت في هذه القصيدة الميسّرة أن أختصر كتاب «التيسير» للداني، فأثمرت بعون الله ما كنت أرجوه منه.

\begin{tanbih}[فائدة]
لهذا يقال: الشاطبية نَظْمُ التيسير. فأصل مادتها من التيسير، مع زيادات قليلة تسمى «زيادات القصيد».
\end{tanbih}

\bayt{69}{وَأَلْفَافُهَا زَادَتْ بِنَشْرِ فَوَائِدٍ}{فَلَفَّتْ حَيَاءً وَجْهَهَا أَنْ تُفَضَّلَا}
\begin{mufradat}
\kalima{ألفافها}{أشجارها الملتفّة، أي أبياتها الكثيرة.}
\kalima{بنشر}{بكثرة.}
\kalima{لفّت وجهها}{سترته.}
\end{mufradat}
\maana زادت القصيدة على التيسير بفوائد كثيرة (كزيادة أحكام، وإشارات لتعليل، ومخارج الحروف)، ومع ذلك غطّت وجهها حياءً أن تُفضَّل على أصلها، أدبًا من الفرع مع الأصل ومن التلميذ مع شيخه.

\bayt{70}{وَسَمَّيْتُهَا حِرْزَ الْأَمَانِي تَيَمُّنًا}{وَوَجْهَ التَّهَانِي فَاهْنِهِ مُتَقَبَّلَا}
\begin{mufradat}
\kalima{الحِرز}{ما يُحفظ فيه الشيء.}
\kalima{الأماني}{جمع أُمنية، وهي ما يُتمنى.}
\kalima{وجه التهاني}{أحسن التهاني (وجه الشيء: أحسنه).}
\kalima{تيمّنًا}{تبرّكًا.}
\kalima{فاهْنِه}{فأعطِه (أصله: فاهنأه).}
\end{mufradat}
\maana سميت القصيدة «\textbf{حِرز الأماني ووجه التهاني}» تبركًا وتفاؤلًا بأن تتحقق فيها أماني طلاب هذا العلم. فيا أيها الطالب أعطها عنايتك متقبلًا لها.

\section{الدعاء}

\bayt{71}{وَنَادَيْتُ اللَّهُمَّ يَا خَيْرَ سَامِعٍ}{أَعِذْنِي مِنَ التَّسْمِيعِ قَوْلًا وَمِفْعَلَا}
\begin{mufradat}
\kalima{أعِذني}{أجِرني واحفظني.}
\kalima{التسميع}{أن يعمل الخير ليسمعه الناس فيمدحوه (الرياء).}
\kalima{مِفْعلًا}{فعلًا.}
\end{mufradat}
\maana يا خير مجيب للدعاء، احفظني من الرياء وطلب السمعة في القول والفعل. وخاف الناظم لمّا أشاد بقصيدته أن يدخله شيء من الرياء فاستعاذ منه.

\bayt{72}{إِلَيْكَ يَدِي مِنْكَ الْأَيَادِي تَمُدُّهَا}{أَجِرْنِي فَلَا أَجْرِي بِجَوْرٍ فَأَخْطَلَا}
\begin{mufradat}
\kalima{الأيادي}{النِّعم.}
\kalima{أجِرني}{احفظني.}
\kalima{الجَور}{الميل عن الحق.}
\kalima{أخطلا}{أقع في الكلام الفاسد.}
\end{mufradat}
\maana يا رب إليك أمدّ يدي، ونِعمُك المتتابعة عليّ هي التي جرّأتني على مدّها. احفظني من الميل عن الحق فأقع في فساد القول.

\bayt{73}{أَمِينَ وَأَمْنًا لِلْأَمِينِ بِسِرِّهَا}{وَإِنْ عَثَرَتْ فَهْوَ الْأَمُونُ تَحَمُّلَا}
\begin{mufradat}
\kalima{أمِين}{لغة في «آمين»: اللهم استجب.}
\kalima{الأمين}{الحافظ الموثوق به.}
\kalima{عثرت}{سقطت وأخطأت.}
\kalima{الأمون}{الناقة القوية التي يؤمَن عِثارها.}
\end{mufradat}
\maana اللهم استجب، وهَبْ أمانًا لمن حفظ أسرار هذه القصيدة ونشرها. وإن وقع فيها خطأ فليتحمّله هذا الأمين ويستره، كما تتحمل الناقة القوية الأثقال.

\bayt{74}{أَقُولُ لِحُرٍّ وَالْمُرُوءَةُ مَرْؤُهَا}{لِإِخْوَتِهِ الْمِرْآةُ ذُو النُّورِ مِكْحَلَا}
\begin{mufradat}
\kalima{الحُرّ}{من لم يستعبده هواه.}
\kalima{المروءة}{كمال الأخلاق.}
\kalima{مرؤها}{صاحبها.}
\kalima{ذو النور}{صاحب الإيمان.}
\kalima{مِكحلًا}{الميل الذي يُكتحل به.}
\end{mufradat}
\maana أقول للكريم: صاحب المروءة لإخوانه كالمرآة يريهم عيوبهم ليصلحوها، وبنور إيمانه يشفي عيوبهم كما يشفي الكحلُ العينَ. وفيه إشارة لحديث: «المؤمن مرآة أخيه» (رواه أبو داود).

\section{وصية القارئ لقصيدته}

\bayt{75}{أَخِي أَيُّهَا الْمُجْتَازُ نَظْمِي بِبَابِهِ}{يُنَادَى عَلَيْهِ كَاسِدَ السُّوقِ أَجْمِلَا}
\bayt{76}{وَظُنَّ بِهِ خَيْرًا وَسَامِحْ نَسِيجَهُ}{بِالْإِغْضَاءِ وَالْحُسْنَى وَإِنْ كَانَ هَلْهَلَا}
\bayt{77}{وَسَلِّمْ لِإِحْدَى الْحُسْنَيَيْنِ إِصَابَةٌ}{وَالُاخْرَى اجْتِهَادٌ رَامَ صَوْبًا فَأَمْحَلَا}
\bayt{78}{وَإِنْ كَانَ خَرْقٌ فَادَّرِكْهُ بِفَضْلَةٍ}{مِنَ الْحِلْمِ وَلْيُصْلِحْهُ مَنْ جَادَ مِقْوَلَا}
\begin{mufradat}
\kalima{المجتاز}{المارّ.}
\kalima{يُنادى عليه}{يُعرض للبيع.}
\kalima{كاسد السوق}{غير مرغوب فيه.}
\kalima{أجمِلا}{قل القول الجميل.}
\kalima{نسيجه}{نظمه (شبّهه بالثوب المنسوج).}
\kalima{الإغضاء}{غضّ الطرف عن العيب.}
\kalima{هلهلًا}{رقيقًا ضعيف النسج.}
\kalima{إصابة}{وصول إلى الصواب.}
\kalima{رام صوبًا فأمحلا}{طلب المطر فأصابه الجدب، أي اجتهد فأخطأ.}
\kalima{خَرْق}{عيب.}
\kalima{ادّرِكه}{تداركه وأصلحه.}
\kalima{من جاد مِقْولًا}{من كان فصيح اللسان متمكنًا من العلم.}
\end{mufradat}
\maana[(الأبيات ٧٥–٧٨)] يا أخي المارّ بقصيدتي وهي تُعرض كأنها بضاعة كاسدة: قل فيها قولًا جميلًا، وأحسن الظن بناظمها، وتجاوز عن عيوبها ولو كان نظمها ضعيفًا كالثوب الرقيق.
وسلّم لي بإحدى الحُسنيين: إما أني أصبت فلي أجران، وإما أني اجتهدت فأخطأت فلي أجر، كما في الحديث: «إذا حكم الحاكم فاجتهد ثم أصاب فله أجران، وإذا حكم فاجتهد ثم أخطأ فله أجر» (متفق عليه).
وإن وجدت عيبًا فتداركه بحلمك، وليصلحه العالِم الفصيح. (وهذا تواضع من الناظم، وإلا فنظمه آية في القوة والإتقان.)

\section{وصايا وآداب}

\bayt{79}{وَقُلْ صَادِقًا لَوْلَا الْوِئَامُ وَرُوحُهُ}{لَطَاحَ الْأَنَامُ الْكُلُّ فِي الْخُلْفِ وَالْقِلَى}
\begin{mufradat}
\kalima{الوئام}{الوفاق والمحبة.}
\kalima{طاح}{هلك.}
\kalima{الأنام}{الناس.}
\kalima{الخُلف}{الاختلاف.}
\kalima{القِلى}{البغض.}
\end{mufradat}
\maana قل صادقًا: لولا الوفاق والمحبة بين الناس لهلكوا جميعًا في الخلاف والبغض. فالوفاق سبب الحياة الطيبة، والخلاف سبب الهلاك.

\bayt{80}{وَعِشْ سَالِمًا صَدْرًا وَعَنْ غِيبَةٍ فَغِبْ}{تُحَضَّرْ حِظَارَ الْقُدْسِ أَنْقَى مُغَسَّلَا}
\begin{mufradat}
\kalima{غِيبة}{ذكر المرء أخاه بما يكره.}
\kalima{فغِب}{ابتعد عن مجالسها.}
\kalima{حِظار القدس}{الجنة (الحِظار: ما يُحاط به، والقدس: الطهر).}
\kalima{أنقى مغسَّلا}{نظيفًا مطهرًا من الذنوب.}
\end{mufradat}
\maana عش سليم القلب من الغش والحقد، وابتعد عن مجالس الغيبة، يُدخلك الله الجنة طاهرًا من الذنوب.

\bayt{81}{وَهَذَا زَمَانُ الصَّبْرِ مَنْ لَكَ بِالَّتِي}{كَقَبْضٍ عَلَى جَمْرٍ فَتَنْجُو مِنَ الْبَلَا}
\maana زماننا زمان الصبر، فمن يأتيك بحالة الصبر التي هي كالقبض على الجمر فتنجو بها من العذاب؟ وفيه إشارة لحديث: «يأتي على الناس زمان الصابر فيهم على دينه كالقابض على الجمر» (رواه الترمذي).

\bayt{82}{وَلَوْ أَنَّ عَيْنًا سَاعَدَتْ لَتَوَكَّفَتْ}{سَحَائِبُهَا بِالدَّمْعِ دِيمًا وَهُطَّلَا}
\bayt{83}{وَلَكِنَّهَا عَنْ قَسْوَةِ الْقَلْبِ قَحْطُهَا}{فَيَا ضَيْعَةَ الْأَعْمَارِ تَمْشِي سَبَهْلَلَا}
\begin{mufradat}
\kalima{ساعدت}{أعانت صاحبها على البكاء.}
\kalima{توكّفت}{قطرت.}
\kalima{سحائبها}{مدامعها (شبّهها بالسحاب).}
\kalima{دِيَمًا}{جمع دِيمة: المطر الدائم.}
\kalima{هُطّلا}{جمع هاطل: المطر المتتابع.}
\kalima{قحطها}{جفافها.}
\kalima{سبهللًا}{فارغة بلا فائدة.}
\end{mufradat}
\maana[(البيتان ٨٢–٨٣)] لو أعانت العينُ صاحبها على البكاء على تقصيره في طاعة الله لسالت دموعها دائمة متتابعة كالمطر، ولكن جفاف العين سببه قسوة القلب من الغفلة. فيا لضياع الأعمار التي تمر فارغة بلا عمل نافع!

\bayt{84}{بِنَفْسِي مَنِ اسْتَهْدَى إِلَى اللهِ وَحْدَهُ}{وَكَانَ لَهُ الْقُرْآنُ شِرْبًا وَمَغْسَلَا}
\bayt{85}{وَطَابَتْ عَلَيْهِ أَرْضُهُ فَتَفَتَّقَتْ}{بِكُلِّ عَبِيرٍ حِينَ أَصْبَحَ مُخْضَلَا}
\begin{mufradat}
\kalima{بنفسي}{أفدي بنفسي.}
\kalima{استهدى}{طلب الهداية.}
\kalima{شِربًا}{نصيبًا من الماء يشربه.}
\kalima{مَغسَلًا}{مكانًا يغتسل فيه، أي يتطهر به من الذنوب.}
\kalima{تفتّقت}{انشقّت.}
\kalima{عَبِير}{أخلاط من الطيب.}
\kalima{مُخضَلًّا}{مبتلًّا.}
\end{mufradat}
\maana[(البيتان ٨٤–٨٥)] أفدي بنفسي من طلب الهداية من الله وحده، وكان القرآن له شرابًا يرويه وماءً يطهره من الذنوب. وطابت له الأرض التي يمشي عليها فانتشر منها الطيب (كناية عن ثناء الناس عليه وكثرة الخير ببركته) حين ابتلّ بنعم الله.

\bayt{86}{فَطُوبَى لَهُ وَالشَّوْقُ يَبْعَثُ هَمَّهُ}{وَزَنْدُ الْأَسَى يَهْتَاجُ فِي الْقَلْبِ مُشْعِلَا}
\bayt{87}{هُوَ الْمُجْتَبَى يَغْدُو عَلَى النَّاسِ كُلِّهِمْ}{قَرِيبًا غَرِيبًا مُسْتَمَالًا مُؤَمَّلَا}
\begin{mufradat}
\kalima{طوبى}{الحالة الطيبة، أو شجرة في الجنة.}
\kalima{همّه}{قصده وإرادته.}
\kalima{زند الأسى}{ما يقدح نار الحزن.}
\kalima{يهتاج}{يثور ويلتهب.}
\kalima{المجتبى}{المختار.}
\kalima{يغدو}{يمرّ ويصبح.}
\kalima{مستمالًا}{يميل الناس إليه.}
\kalima{مؤمَّلًا}{يُرجى عند الشدائد.}
\end{mufradat}
\maana[(البيتان ٨٦–٨٧)] فطوبى لهذا المستهدي، يبعث الشوقُ إرادتَه إلى ما أعدّه الله لأهل طاعته، ويشتعل قلبه حزنًا على ما ضاع من عمره. هو المختار عند الله، يعيش بين الناس قريبًا منهم بتواضعه، غريبًا بندرة حاله، يحبه الناس ويرجون دعاءه.

\bayt{88}{يَعُدُّ جَمِيعَ النَّاسِ مَوْلًى لِأَنَّهُمْ}{عَلَى مَا قَضَاهُ اللهُ يُجْرُونَ أَفْعُلَا}
\bayt{89}{يَرَى نَفْسَهُ بِالذَّمِّ أَوْلَى لِأَنَّهَا}{عَلَى الْمَجْدِ لَمْ تَلْعَقْ مِنَ الصَّبْرِ وَالْأَلَا}
\begin{mufradat}
\kalima{مولى}{سيدًا (على المعنى الأول)، أو عبدًا مملوكًا لله (على المعنى الثاني).}
\kalima{أفعُلا}{أفعالًا.}
\kalima{لم تلعق}{لم تذق.}
\kalima{الصَّبِر}{عصارة شجر مُرّ.}
\kalima{الأَلا}{شجر حسن المنظر مُرّ الطعم.}
\end{mufradat}
\maana[(البيتان ٨٨–٨٩)] هذا المختار يعدّ الناس كلهم سادة تواضعًا فلا يحتقر أحدًا، أو يعدّهم عبيدًا لله لا يملكون له نفعًا ولا ضرًّا فلا يخافهم ولا يتملقهم، لأن أفعالهم تجري بقضاء الله.
ويرى نفسه أولى باللوم من غيره؛ لأنها لم تتحمل مرارة الصبر في طلب المجد عند الله، فيشتغل بعيب نفسه عن عيوب الناس.

\bayt{90}{وَقَدْ قِيلَ كُنْ كَالْكَلْبِ يُقْصِيهِ أَهْلُهُ}{وَمَا يَأْتَلِي فِي نُصْحِهِمْ مُتَبَذِّلَا}
\begin{mufradat}
\kalima{يُقصيه}{يُبعده.}
\kalima{ما يأتلي}{ما يقصّر.}
\kalima{متبذّلًا}{باذلًا كل جهده.}
\end{mufradat}
\maana قيل في المثل: كن كالكلب في الوفاء؛ يطرده أهله ويؤذونه، وهو لا يقصّر في حراستهم ونصحهم. والمقصود: الحث على بذل الجهد في طاعة الله والصبر على البلاء، والإخلاص للناس ولو أساؤوا.

\bayt{91}{لَعَلَّ إِلَهَ الْعَرْشِ يَا إِخْوَتِي يَقِي}{جَمَاعَتَنَا كُلَّ الْمَكَارِهِ هُوَّلَا}
\bayt{92}{وَيَجْعَلُنَا مِمَّنْ يَكُونُ كِتَابُهُ}{شَفِيعًا لَهُمْ إِذْ مَا نَسُوهُ فَيَمْحَلَا}
\begin{mufradat}
\kalima{يقي}{يحفظ.}
\kalima{هُوَّلًا}{جمع هائل: المخيف المفزع.}
\kalima{يَمْحَل}{يشكو ويسعى بصاحبه إلى الله (محَل به: وشى به).}
\end{mufradat}
\maana[(البيتان ٩١–٩٢)] أرجو يا إخوتي أن يحفظنا الله من كل المكاره المفزعة، وأن يجعلنا ممن يكون القرآن شفيعًا لهم يوم القيامة لأنهم لم ينسوه، لا ممن يشكوهم القرآن إلى ربه. وفيه إشارة لحديث: «القرآن شافع مشفَّع، وماحِلٌ مصدَّق».

\bayt{93}{وَبِاللهِ حَوْلِي وَاعْتِصَامِي وَقُوَّتِي}{وَمَا لِيَ إِلَّا سِتْرُهُ مُتَجَلِّلَا}
\bayt{94}{فَيَا رَبِّ أَنْتَ اللهُ حَسْبِي وَعُدَّتِي}{عَلَيْكَ اعْتِمَادِي ضَارِعًا مُتَوَكِّلَا}
\begin{mufradat}
\kalima{حولي}{تحوّلي من حال إلى حال (من المعصية إلى الطاعة).}
\kalima{اعتصامي}{امتناعي من كل ما يَشين.}
\kalima{متجلّلًا}{متغطّيًا.}
\kalima{حسبي}{كافيَّ.}
\kalima{عُدّتي}{ما أُعِدّه لدفع الشدائد.}
\kalima{ضارعًا}{متذللًا.}
\end{mufradat}
\maana[(البيتان ٩٣–٩٤)] تحوّلي من المعصية إلى الطاعة، وامتناعي من الشر، وقوتي على الخير — كل ذلك بالله وحده (وهذا معنى «لا حول ولا قوة إلا بالله»). وليس لي ما أعتمد عليه إلا ستره الذي سترني به في الدنيا، وأرجوه في الآخرة.
يا رب أنت كافيّ في كل أمر، وعُدّتي في كل شدة، وعليك وحدك اعتمادي، متذللًا متوكلًا عليك.

\begin{khulasa}[ختام المقدمة]
بهذا تنتهي مقدمة الشاطبية (٩٤ بيتًا). وأهم ما يجب ضبطه منها قبل الدخول في الأبواب:
\begin{enumerate}
\item أسماء القراء السبعة ورواتهم بالترتيب (الأبيات ٢٥–٤٠).
\item رموز «أبجد» ورموز الجماعة (الأبيات ٤٥–٥٥) — راجع جدول دليل الاستخدام.
\item قواعد الأضداد والإطلاق والتقييد (الأبيات ٥٧–٦٣).
\end{enumerate}
\end{khulasa}
